\documentclass[12pt]{article}
\usepackage[utf8]{inputenc}
\usepackage[T1]{fontenc}
\usepackage[margin=1in]{geometry}
\usepackage{setspace}
\usepackage{xcolor}
\usepackage[colorlinks=true, urlcolor=blue!60!black, linkcolor=blue!60!black]{hyperref}
\setlength{\parindent}{0em}
\setlength{\parskip}{0.6em}
\renewcommand{\baselinestretch}{1.0}
\begin{document}
\begin{center}
{\LARGE \textit{Wisconsin v. Yoder}, 406 U.S. 205 (1972)\par}
\vspace{0.5em}
{\large Case Brief --- Two-Page Edition\par}
\end{center}

\section{Memory Jogger}
Applying strict scrutiny under the Free Exercise Clause, the Court held that Wisconsin could not punish Amish parents for keeping their children out of high school after eighth grade (\href{https://supreme.justia.com/cases/federal/us/406/205/}{406 U.S. 205}).

\section{Detailed Case Facts}
Amish parents Jonas Yoder, Wallace Miller, and Adin Yutzy withdrew their 14- and 15-year-olds after eighth grade, violating Wisconsin's attendance-until-16 law. Uncontradicted experts testified that high school would pull adolescents from the community before adult baptism and threaten Amish survival, and the State conceded the beliefs were sincere.

\section{Procedural History}
The parents were convicted and fined \$5 each; the Wisconsin Supreme Court reversed, and the U.S. Supreme Court affirmed.

\section{Judicial Votes}
Affirmed, 6--1. Burger wrote for the Court; Stewart and White concurred; Douglas dissented in part; Powell and Rehnquist did not participate.

\section{Holding}
Only ``interests of the highest order and those not otherwise served'' can overbalance Free Exercise claims, and Wisconsin did not show how exempting the Amish from one or two more years of school would harm its goals.

\section{Analysis of Opinions}
Burger linked free exercise to the parental rights recognized in \textit{Pierce}. A secular way of life like Thoreau's ``does not rise to the demands of the Religion Clauses,'' but the Amish life is ``one of deep religious conviction.'' Belief and action ``cannot be neatly confined in logic-tight compartments,'' and facial neutrality is no defense. The Amish were law-abiding and self-sufficient, the State's fears ``highly speculative,'' and \textit{Prince} did not reach this record. Stewart stressed that only the parents were prosecuted; White called the case ``close.'' Douglas would hear the children: ``[i]t is the future of the student, not the future of the parents, that is imperiled.'' \textit{Mahmoud v. Taylor} (2025) later revitalized \textit{Yoder}.

\section{Examples: Future Applications}
\textbf{Same-Side:} prosecuting Amish parents under a new law requiring attendance until age 18 fails. \textbf{Opposite-Side:} a secular commune inspired by Thoreau's claim fails for lack of religious belief. \textbf{Fence-Sitter:} a school measles-vaccine requirement that allows medical but not religious exemptions, which pits \textit{Prince v. Massachusetts} against \textit{Fulton v. City of Philadelphia}.

\section{Critique}
Relying on the Amish pedigree and crime rate risks disfavoring newer faiths, and excluding secular conscience strains Establishment Clause neutrality. Deferring to parents over children, and \textit{Mahmoud}'s expansion, raise limiting-principle questions.

\section{Key Quotations}
\begin{quote}
``The essence of all that has been said and written on the subject is that only those interests of the highest order and those not otherwise served can overbalance legitimate claims to the free exercise of religion.''
\end{quote}
\noindent --- Chief Justice Burger, Opinion of the Court, 406 U.S. at 215.

\end{document}
